\subsection{Commercial Recommendation Systems}

Large-scale commercial platforms have driven much of the practical advancement in recommendation system design.

Amazon's recommendation engine~\cite{linden2003amazon} pioneered item-to-item collaborative filtering at scale, generating ``customers who bought X also bought Y'' associations from co-purchase logs.
The central insight, that item-to-item similarity is far more scalable to compute offline than user-to-user similarity, directly motivates the Cleora co-purchase graph used in Pipeline~A of the proposed system.
Subsequent iterations of Amazon's system incorporate session-based signals, click-through prediction models, and contextual bandit approaches for online exploration, but the co-purchase neighbourhood structure remains a core retrieval mechanism.

Spotify's music recommendation system~\cite{van2013deep} combined collaborative filtering with audio content features derived from a convolutional neural network, demonstrating that modality-specific feature extraction could improve cold-start performance for new tracks with no interaction history.
This multimodal motivation aligns with the use of BGE-M3 textual embeddings and CLIP visual embeddings in the proposed system: items with zero interactions can still be retrieved by content similarity.

Netflix's recommendation research~\cite{gomez2016deep} has emphasised personalised ranking at scale, deploying deep learning models trained on hundreds of millions of user-item interaction pairs.
The prize-winning matrix factorisation approaches from the Netflix Prize competition established the importance of latent factor models for capturing long-term preference patterns.
While the proposed system does not use a dedicated user-embedding matrix, the Cleora co-purchase graph implicitly encodes similar co-occurrence structure within its item embeddings.

Goodreads, as a domain-specific book discovery platform, presents challenges specific to the book domain that motivate this project: extremely long catalogue tails (millions of titles), significant multilingual demand, and user behaviour driven more by genre-level affinity than by fine-grained interaction frequency.
These domain characteristics favour the content-anchored retrieval strategy of Pipeline~A and the multilingual BGE-M3 encoder over purely behavioural collaborative filtering.

\subsection{Multimodal Recommendation in the Literature}

The incorporation of multiple content modalities (text, image, and structured attributes) into recommendation systems has received substantial attention following the success of vision-language models.

\paragraph{Vision-language embedding for items.}
Radford et al.~\cite{radford2021clip} demonstrated that contrastive pre-training on image-text pairs produces aligned visual and textual embedding spaces, enabling zero-shot image-to-text retrieval.
Applied to product recommendation, this alignment allows item cover images and textual descriptions to be compared in a shared space.
The CLIP index in the proposed system exploits this property: a user can submit a book cover photograph and retrieve semantically related titles even when the exact item is not in the catalogue.

\paragraph{Language-model-augmented recommendation.}
The BLaIR encoder~\cite{hou2024blair} and its successors applied transformer-based language models to encode product descriptions, review text, and titles into item embeddings suitable for dense retrieval.
The transition from BLaIR to BGE-M3 in the proposed system reflects a broader trend towards multilingual, multi-granularity encoders: BGE-M3~\cite{chen2024bge} supports 100+ languages and achieves superior performance on both dense retrieval and cross-lingual semantic similarity tasks, addressing the Vietnamese-query failure mode of the legacy system.

\paragraph{Graph neural networks for collaborative signals.}
LightGCN~\cite{he2020lightgcn} showed that a simplified graph convolutional network propagating signals over the user-item bipartite graph could outperform more complex graph neural network variants on standard collaborative filtering benchmarks.
Cleora~\cite{rychalska2021cleora} achieves comparable representational quality through a simpler iterative normalisation scheme that is orders of magnitude faster to train on large graphs, making it tractable for the 375{,}280-item co-purchase graph used in Pipeline~A without requiring a full GNN training loop.

\paragraph{Hybrid retrieval with re-ranking.}
Systems combining sparse lexical matching with dense semantic retrieval have become standard practice in information retrieval.
Cormack et al.~\cite{cormack2009rrf} proposed Reciprocal Rank Fusion (RRF) as a lightweight, parameter-free method for combining result lists from multiple retrievers.
The search pipeline in the proposed system applies RRF to fuse HNSW semantic search results with Tantivy BM25 results, reproducing the retrieval architecture used in industrial-scale neural search systems such as Elastic Learned Sparse Retrieval and OpenAI's embedding-based retrieval pipelines.
